\begin{pastexam}[label={px:ust-115-ee-ds-40}]{臺聯大電機類 115 學年度 第 40 題}
（單選；2 分）

The following piece of code is a version of \textbf{Insertion Sort Algorithm} that used a linked
list to store the sorted data. Fill in the two blank areas on Lines 14 and 20 to make the
\texttt{insert()} function maintain the list as a sorted list in \textbf{ascending} order. Nodes
with equal values (\texttt{val}) must be ordered according to their insertion time (the earlier
ones are placed before the later ones).
\begin{minted}[linenos, numbersep=1.5mm, xleftmargin=2em, fontsize=\small]{c}
#include <stdio.h>
#include <stdlib.h>
struct Node {
    int val;
    struct Node* next;
};
struct Node* createNode(int x) {
    struct Node* node = (struct Node*)malloc(sizeof(struct Node));
    node->val = x;
    node->next = NULL;
    return node;
}
struct Node* insert(struct Node* head,struct Node* item){
    if(head==NULL || head->val ____ item->val){
        item->next=head;
        head=item;
    }
    else {
        struct Node* curr=head;
        while(curr-next != NULL && curr->next->val ____ item->val) {
            curr=curr->next;
        }
        item->next=curr->next;
        curr->next=item;
    }
    return head;
}
\end{minted}
\begin{enumerate}[label=(\Alph*)]
\item \texttt{Line 14:\underline{ > }, Line 20:\underline{ <= }}
\item \texttt{Line 14:\underline{ > }, Line 20:\underline{ < }}
\item \texttt{Line 14:\underline{ < }, Line 20:\underline{ <= }}
\item \texttt{Line 14:\underline{ < }, Line 20:\underline{ >= }}
\item \texttt{Line 14:\underline{ >= }, Line 20:\underline{ < }}
\end{enumerate}
\end{pastexam}

\begin{pastsoln}
答案是 (A)。

原卷第 20 行的 \texttt{curr-next} 是 \texttt{curr->next} 的筆誤，下面照 \texttt{curr->next} 解。

這是在排好的串列裡插入一個節點（資結書虛擬碼~5.4）。 要求有兩個：由小到大；相同的值，先插入的
排前面（也就是插入排序要穩定）。 所以新節點 \texttt{item} 要放在「所有 $\le$ 它的節點之後、第一個
$>$ 它的節點之前」。
\begin{itemize}
\item \textbf{第 14 行}：什麼時候 \texttt{item} 要放在最前面？ 串列是空的，或頭節點\textbf{嚴格大於}
\texttt{item}。 如果頭節點和 \texttt{item} 相等，頭節點比較早插入，\texttt{item} 必須排在它後面，
不能放到最前面。 所以填 \texttt{>}。
\item \textbf{第 20 行}：\texttt{curr} 往後走的條件是下一個節點還應該排在 \texttt{item} 前面，也就是
\texttt{curr->next->val <= item->val}（相等的舊節點也要跳過）。 迴圈停下時，\texttt{curr->next}
是第一個 $>$ \texttt{item} 的節點（或 \NIL），接在 \texttt{curr} 後面正確。 所以填 \texttt{<=}。
\end{itemize}

用依序插入 $3, 1, 3', 2, 1'$（加撇號的是後插入的相同值）檢查各選項，正確結果應為
$1, 1', 2, 3, 3'$：
\begin{itemize}
\item (A) \texttt{>}、\texttt{<=}：得 $1, 1', 2, 3, 3'$。 對。
\item (B) \texttt{>}、\texttt{<}：第 20 行遇到相等的就停，新節點插在相等的舊節點前面。 得
$1, 1', 2, 3', 3$，$3'$ 跑到 $3$ 前面，不穩定。 （$1'$ 沒出問題，是因為它在第 14 行就被判定不放最前面、
且 $1$ 後面是 $2$；一般情形不保證。） 錯。
\item (C) \texttt{<}、\texttt{<=}：插入 $1$ 時頭節點 $3 < 1$ 不成立，進到 \texttt{else}，$1$ 被接在
$3$ 後面，得 $3, 1, \ldots$，連大小順序都錯。 錯。
\item (D) \texttt{<}、\texttt{>=}：這是由大到小的寫法，得 $3, 3', 2, 1, 1'$。 錯。
\item (E) \texttt{>=}、\texttt{<}：相等時放到最前面或停在相等節點之前，後插入的相同值都跑到前面，
得 $1', 1, 2, 3', 3$。 錯。
\end{itemize}
\end{pastsoln}
